% ============================================================
%  开题报告正文骨架（AAD 毕设）
%
%  本文件当前是**框架**，不是成稿：
%    · 四节结构与两级标题完全照官方 2027 届模板，不动；
%    · 每个小节下用 \tjformslot 标出「这一段要写什么」，灰色、【待写】开头；
%    · 技术路线图用 \tjformfigslot 留图位，还没画；
%    · 参考文献只留位置，真实条目等检索完再填。
%
%  🔴 成稿前必做两件清理：
%     grep -n 'tjformslot\|tjformfigslot' chapters/proposal_body.tex   → 应为 0
%     grep -n 'tjformnote' chapters/proposal_body.tex                  → 应为 0（模板说明灰字，成稿删除）
%
%  🔴 防泄露（导师 2026-11 微信要求）：本文件不得写入 AAD 论文未公开的核心创新点——
%     消息可达性判据的具体定义、H1–H4 因果识别框架、十四类消息生命周期分类、
%     区域风险聚合的具体公式、以及现有量化结论。只写研究问题、技术路线、公开指标。
% ============================================================

\tjformsection{一、\tjheadertext 课题背景}
\tjformnote{根据实际情况写为毕业设计或毕业论文，下同（此项成稿删除）}

\tjformsubsection{1．选题背景与研究意义}
\tjformslot{开篇交代混合交通的必然性：自动驾驶、智能网联与人工智能技术发展，未来道路交通将长期是自动驾驶车辆与人工驾驶车辆共存。}
\tjformslot{再讲问题从哪来：两类车在感知能力、决策方式、响应时间、驾驶行为上差异明显，局部风险不再局限于单车或单一冲突，可能通过跟驰、制动、换道与交互行为向其他车辆和局部区域传播，并在特定条件下放大成级联响应。}
\tjformslot{然后指出缺口：只做单车风险独立分析，刻画不了整体安全演化；所以要研究局部风险如何产生、传播、累积，以及交通状态、车辆类型、交互关系对传播的影响。}
\tjformslot{最后落到意义：理解混合交通的安全机理、提高系统韧性，并为安全分析与运行管理提供理论与实验依据。}

\tjformsubsection{2．国内外研究现状（文献综述）}
\tjformslot{本节按四条线索组织，每条讲清「做了什么 — 局限在哪 — 本研究因此补什么」。真实文献检索完成后逐条填充，不得凭印象引用。}
\tjformslot{线索一 混合交通风险评价：替代安全指标（TTC、DRAC 等）的提出、适用条件与局限。}
\tjformslot{线索二 交通扰动与风险传播：扰动如何在车流中传播、现有研究做到什么尺度、缺什么。}
\tjformslot{线索三 微观交通仿真：SUMO 等平台在混合交通场景构建与校核中的用法与不足。}
\tjformslot{线索四 面向风险抑制的控制：现有干预思路与评价维度（安全性、效率、平稳性）。}
\tjformslot{小结：指出现有研究在「风险跨车辆、跨区域传播的可测与可干预」上的空白，由此引出本课题。}

\tjformsection{二、\tjheadertext 方案介绍}

\tjformsubsection{1．主要研究内容}
\tjformslot{与任务书「二、毕业论文应完成的主要内容」五条一一对应，每条一两句写清做什么、产出什么。}
\tjformslot{内容一 研究综述：调研混合交通、交通扰动传播、车辆交互风险、替代安全指标与交通仿真，梳理方法及适用范围，明确研究对象与技术路线。}
\tjformslot{内容二 风险测量方法：基于真实轨迹整理速度、加速度、车间距离与交互关系，选适用于典型跟驰冲突场景的安全评价指标，研究风险事件识别方法，给出指标合理取值范围。}
\tjformslot{内容三 仿真环境构建：基于 SUMO 等平台构建含自动驾驶与人工驾驶车辆的典型场景，设置车辆组成、交通流状态、驾驶行为与局部扰动，并用真实轨迹校核关键运动参数与风险状态。}
\tjformslot{内容四 传播规律与影响因素：设计不同交通状态、车辆组成与局部扰动条件下的实验，记录风险时空演化，分析传播、累积、衰减特征与影响范围、速度、强度。}
\tjformslot{内容五 抑制方法探索：依据传播规律研究利用车辆状态与风险信息提前干预的控制思路，对比安全性、效率与运行平稳性，评价抑制效果。}

\tjformsubsection{2．研究方案}
\tjformslot{本小节按「数据 — 方法 — 平台 — 实验 — 控制 — 路线图 — 可行性」七块展开。}
\tjformslot{数据：真实车辆轨迹数据（来源、字段、覆盖场景、预处理与异常缺失处理口径）。}
\tjformslot{方法：风险指标计算与风险事件识别；指标取值范围标定；不写论文未公开的判据定义。}
\tjformslot{平台：SUMO 版本、路网与场景构建方式、车辆行为参数设置、与真实轨迹的校核方法。}
\tjformslot{实验设计：交通状态、车辆组成、局部扰动三类因子如何组合；重复性与随机种子口径；评价指标与统计方法。}
\tjformslot{控制方法：干预信号从哪来、如何作用到车辆、对照设置与评价维度。}
\tjformfigslot{技术路线图（数据 → 指标 → 仿真 → 实验 → 控制，用代码画，图内不写核心算法细节）}
\tjformslot{可行性分析：数据已有、平台已有、方法有前期基础、时间安排可行（对应任务书进度）。}

\tjformsubsection{3．工作进度安排}
\tjformslot{照任务书「四、毕业论文的进度安排」七段填写，可整理成表；时间区间必须与任务书一致，不改。}

\tjformsection{三、\tjheadertext 主要参考文献}
\tjformrefnote{正文首行缩进2字符。参考文献悬挂0.74厘米，参考文献格式请参考GB/T 7714-2025《信息与文献 参考文献著录规则》，以下为样例。成稿请删除。}
\begin{tjformrefs}
\tjformslot{真实文献检索完成后填入，按 GB/T 7714-2025 著录，预计 15–20 条；每条必须能点开验证（DOI / arXiv 编号 / 期刊卷期页）。}
\end{tjformrefs}

\tjformsection{四、审核意见}

\begin{tjreviewbox}
  \tjformlabel{指导教师意见：}\tjformhint{（课题难度是否适中、工作量是否饱满、进度安排是否合理、工作条件是否具备、是否同意开题等）}

  \vspace{8\baselineskip}
  \tjformsignature{指导教师签名}{8em}
\end{tjreviewbox}

\begin{tjreviewbox}
  \tjformlabel{专业审核意见：}\tjformhint{（适当对各项内容给与意见和评价，字数不限。）}

  \vspace{6\baselineskip}
  \tjformsignature{负责人签名}{8em}
\end{tjreviewbox}
